\section{Completion of the current DC risk bound}
\begin{proof}[Proof of the regularized DC risk bound below a fixed logarithmic scale]
Fix any constant $A_0<\infty$ and first work on
$d\le A_0\log n$, still with $d\to\infty$.
Let $\mathcal E$ collect the aggregate geometric conditions of
\cref{prop:sparse-floor}, their recomputed versions after deletion of
at most three vertices, two-sided total degree, and the true block-total event used
in the fitting lemmas. None imposes a simultaneous lower bound on all
raw degrees or a simultaneous incoherence bound. For every prescribed
fixed $r$, its failure probability is $O(n^{-r})$ after increasing the
fixed constants. For the block totals this also follows directly from
the independent-edge Bernstein bound at a sufficiently large constant
times $n$: its exponent is at least $cn/d$, and $np^2\le C_1$ gives
$n/d\ge d/C_1$; in the present range it is much larger than $\log n$.

The initializer and invariant-basin result in
\cref{restore:fit-invariant} apply deterministically on $\mathcal E$.
For a tested target $i$, impose only its upper-degree bound
$D_i\le Ld$, its local normalized row bound, and the condition that its
incident row hits no exceptional deleted eigenvector row. By
\cref{restore:tail-energy-lemma,restore:tail-hit-corollary,restore:tail-target-corollary}, the average probability cost
of the two row conditions is $C_He^{-Hd}$ for any prescribed fixed $H$,
after choosing their fixed thresholds sufficiently large. For the
LOO hit condition, the exceptional set consists of deleted rows with
$\sqrt n\norm{(U_B)_{j\cdot}}_2>R$; the floor implies
\eqref{restore:loo-hit} on its complement. Its expected hit mass
is bounded by the row-count argument in
\eqref{restore:tail-coefficient-count}.
The upper
degree tail has the same property by the single-row Bernoulli Chernoff
bound. Polynomial factors in $d$ are absorbed by choosing a larger
exponential rate at the preceding step. No union over the $n$ tested
targets is taken: their error probabilities are averaged.

On these local events, the fitted-score lemma
\cref{restore:fit-score} replaces every estimated, possibly
data-dependent contrast by a deterministic true block contrast at
$O(1)$ cost. The local LOO comparison
\cref{restore:loo-corrected} replaces its actual full-graph spectral
action by an independent training action at another $O(1)$ cost.
For the DC profile contrast, the new score is the true raw-row
linear profile contrast plus
\[
 \sum_{j\ne i}A_{ij}e_j+O(1),\qquad
 e=(H_B-I)(w-\bar w\one).
\]
The coefficients are measurable from the deleted graph and independent
of the incident Bernoulli row. For each target use the analogous
training event $\mathcal E_i$ formed only from that deleted graph and
its further deletions. In applying the independent-row bound,
drop the full-graph event after the deterministic error inclusion and
condition only on the training graph on $\mathcal E_i$.
The geometry cost is
\[
 \varepsilon_{\rm geom}
 =\PP(\mathcal E^c)+\frac1n\sum_i\PP(\mathcal E_i^c)
 \le C_rn^{-r}.
\]
This does not condition the Bernoulli row on a full-graph event.
The tilted mean and quadratic cost are
bounded by \eqref{eq:floor-tilt-moments}; here $np^2$ is bounded.
The exceptional coefficient hit-mass hypothesis of
\cref{lem:sparse-truncated-transfer} is now a conclusion of
\cref{restore:tail-hit-corollary}, rather than an unproved premise. That sharp
transfer gives, after averaging the local target exceptions and taking
a fixed-$K$ union over competitors,
\[
 \E\ell(\widehat z,z)
 \le \frac Cn\sum_i\sum_{b\ne z_i}
                  d^{-1/2}e^{-J_{i,z_i\to b}}
       +C_He^{-Hd}+\varepsilon_{\rm geom}.
\]
Here $J_{i,a\to b}$ is the exact directional profile exponent in
\eqref{dcsbm:profile-exponent}. By \cref{dcsbm:profile-tail}, the same
pairwise directional sharp scale is a lower bound, up to a constant,
for the actual profile oracle. Since $K$ is fixed, the
displayed sum is bounded by a constant times its average risk.
Only the averaged target-row tail is used; no pointwise uniform tail
for heterogeneous degree parameters is asserted.

Uniformly on the compact DC class, the reference average risk is at
least $c d^{-1/2}e^{-C_Jd}$. Choose $H>C_J+1$ and then choose fixed
$r>A_0(C_J+1)+2$. For $d\le A_0\log n$,
\[
 \frac{e^{-Hd}+n^{-r}}{d^{-1/2}e^{-C_Jd}}\longrightarrow0.
\]
This closes the unconditional expected-risk estimate throughout the
entire range below any fixed multiple of $\log n$. In particular, no
positive lower coefficient of $\log n$ is required.
\end{proof}

\begin{proof}[Completion over the remaining degree scales]
Fix $A_0>0$ and use the preceding sparse-range argument when
$d\le A_0\log n$. For $d>A_0\log n$, replace its polynomially
accurate aggregate event by $\mathcal E_H$ from
\cref{restore:high-aggregate}. The proof of
\cref{restore:tail-energy-lemma} uses only the deterministic aggregate
geometry on the event; the innovation concentration is applied for
every training graph before restricting to the event. It therefore
holds unchanged with $\mathcal E_H$ and the corresponding
training-measurable events. Its coefficient-hit and target-row
corollaries, the scalar LOO comparison, all fitted-score bounds,
and the invariant basin likewise apply.

The resulting risk upper bound is the appropriate oracle term plus
$C_{H'}e^{-H'd}+C_He^{-Hd}$, where $H'$ can be prescribed by choosing
the fixed local thresholds. On the fixed positive compact DC model
class the oracle risk is at least $cd^{-1/2}e^{-C_Jd}$.
Choosing $H,H'>C_J+1$ pays both remainders. This proves the stated
all-scale profile-oracle bound. The known-target-degree comparison follows
from the separate equal-exposure calculation in
\cref{dcsbm:equal-exposure}. No logarithmic degree lower bound is imposed.
\end{proof}
